\section{Evaluation Protocol}
\label{sec:protocol}

\subsection{Matching criterion: IoU is the wrong question}
\label{sec:criterion}

A detection metric asks whether a predicted box \emph{localizes} an animal; a survey asks
whether the animal was \emph{noticed}. On a 27-pixel target these diverge sharply: a box
offset by six pixels falls below IoU 0.5 and is scored as both a false positive and a false
negative, while for counting it has correctly found the animal. We therefore evaluate under
four matching rules --- IoU$\geq$0.50, IoU$\geq$0.30, centre-in-box and \emph{any overlap} ---
and treat the last as the counting criterion, reporting it as presence or counting recall
and never as mAP, with the strict criterion alongside it everywhere.

The criterion is validated against centre-in-box, an independent counting-oriented rule that
accepts a prediction only when its centre falls inside the animal and so rejects oversized
boxes outright. The two agree to within 0.032 F1 across all twelve models, so the
conclusions rest on two criteria agreeing rather than on one being chosen. Precision is
reported alongside recall throughout, which additionally constrains box quality.

\subsection{Decomposing counting error}
\label{sec:decomposition}

The paper's main instrument. Let $G$ be the ground-truth animals in a video and $C$ the
candidate tracks a pipeline produces. For each $g\in G$ we define:

\begin{itemize}\itemsep2pt \parskip0pt
\item \reachedq --- $g$ is touched, on at least one frame, by at least one candidate. If
an animal is not \reachedq, no downstream stage can recover it; this measures
\textbf{detection and tracking}.
\item \primaryq --- among the candidates touching $g$, the one that covers $g$ on the most
frames is not simultaneously the best cover for some other animal. Two animals sharing a
single best-covering candidate are indistinguishable to any decision rule that emits one
count per track; this measures \textbf{association}.
\item \countedq --- the confirmation stage accepted a candidate covering $g$. This measures
\textbf{confirmation}, and it is the end-to-end result.
\end{itemize}

The three are nested, $\countedq \le \primaryq \le \reachedq \le |G|$, and the gaps
attribute error to the three stages. We report all three for every configuration; reporting
only the last, as is conventional, makes a detection failure and a confirmation failure look
identical.

\subsection{Counting metrics}

We report mean absolute error in animals per video (MAE), signed bias, and a
\emph{counted} fraction. The latter is $\sum_v \min(\hat{c}_v, c_v) / \sum_v c_v$, capped
per video so that over-counting in one video cannot mask misses in another.

We also report a stricter identity-matched variant, the fraction of individual animals
covered by an \emph{accepted} candidate, which \cref{sec:counting} shows is 10 points
lower; the capped aggregate credits a video for predicting the right \emph{number} even
when the accepted tracks do not land on the animals.

\subsection{The held-out protocol}
\label{sec:heldout}

Counting runs process all 32 videos, but the detector trained on 19 of them, and the
hand-tuned confirmation rule is conventionally swept on the same videos it is reported on.
Both leaks inflate the result. Our protocol closes both: \textbf{sweep the rule on the 19
detector-training videos, freeze it, and report on the 13 unseen videos}. Learned
confirmers are trained and thresholded under the identical split, so the comparison in
\cref{sec:counting} is like-for-like. The protocol is worth 23 points: counting accuracy
measured over all 32 videos is 89.0\%, and on the 13 the detector never saw it is 66.3\%.

